\documentclass[a4paper, 11pt]{report}
\usepackage[utf8]{inputenc}
\usepackage[T1]{fontenc}
\usepackage[french]{babel}

\begin{document}

\noindent
UNIVERSITÉ DE LOMÉ \hfill ANNÉE UNIVERSITAIRE \\
FACULTÉ DES SCIENCES \hfill 2021-2022

\begin{center}
  Examen de Phy 112 "Électromagnétisme et Relativité" \\ \vspace{0.2 cm}
  \bf \underline{Partie B: Relativité restreinte} \\ \vspace{0.2 cm}
  \it (Parcours Physique et chimie)
\end{center}

\subsubsection*{\underline{Questions de cours}}

\begin{enumerate}
  \item Énoncer le principe de relativité de Galilée.
  \item Qu'entend on par référentiel de l'Ether ? \\
  Quelle est la vitesse de la lumière dans ce référentiel ?
  \item Quel était l'objectif de l'expérience de Michelson et Morley réalisée en 1887 ? \\
  Quelle conclusion peut-on tirer de cette expérience ?
  \item Quels sont les postulats fondamentaux de la Relativité Restreinte (On dit aussi postulat d'Einstein) ? La transformation de Galilée est-elle en accord avec les postulats d'Einstein ? Justifier votre réponse.
  \item Quelle est la transformation d'espace-temps qui est en accord avec les postulats d'Einstein ?
  \item Donner sans démontrer les relations de la transformation de Lorentz en termes de $ \beta = u/c $ et $ \gamma = \frac{1}{\sqrt{1-{u^2}/{c^2}}} $
\end{enumerate}

\subsubsection*{\underline{Exercices d'application}}

\begin{enumerate}
  \item[A1] Les coordonnées d'un éclair mesurés par $O$ sont $ x = 300 km $ , $ y = 150 km $ , $ z = 100 km $ , à $ t = 4.10^{-3}s $ . \quad {\bf Quelles sont les coordonnées spatio-temporelles} déterminées par un observateur $O'$ , se déplaçant par rapport à $O$ , avec une vitesse de 0,6c , le long de l'ace $x' - x$ commun ?
  \item[A2] Un Observateur $O$  repère deux événements séparés dans l'espace et le temps de $ 20 km $ et $ 6.10^{-5} s $ . \quad {\bf Quelle doit être la vitesse} de déplacement par rapport à $O$ , d'un  observateur $O'$ pour que les deux événements lui paraissent simultanés ?
  \item[A3] 
  \begin{enumerate}
    \item[(a)] Un éclair est observé par $O'$ frappe en $ x' = 60 m $ , $ y' = z' = 0$ , $ t' = 8.10^{-8} s $ . $O'$ se déplace sur l'axe $ x - x' $ (commun avec $O$) avec une vitesse de 0,5c par rapport à $O$ . \quad {\bf Quelles sont les coordonnées spatio-temporelles} de lévénement pour $O$ ?
    \item[(b)] Si un second éclair frappe en $ x' = 10 m $ et $ t' = 2.10^{-7} s $ pour $O'$ , {\bf quel sera l'intervalle de temps entre les deux événements} pour $O$ ?
  \end{enumerate} 
\end{enumerate}

\newpage

\noindent
{\bf UNIVERSITÉ DE LOMÉ \hfill ANNÉE UNIVERSITAIRE \\
FACULTÉ DES SCIENCES \hfill 2020-2021 }

\begin{center}
  {\bf TD de INF203 Algorithmes et Programmation Fortran}
\end{center}

\subsubsection{\underline{Exercice 1}}

\noindent
L'énergie potentielle inermoléculaire peut être calculée par la formule de Renards-jones: 
$$ E(R) = 4\varepsilon((\sigma/R)^2 - (\sigma/R)^6) $$
où $\varepsilon$ et $\sigma$ sont les paramètres caractéristiques des molécules mises en jeu. \\

\vspace{0,5 cm}

\noindent
Écrire un programme permettant de calculer $E(R)$ pour des valeurs de $R$ comprises entre 1 et 5 \AA \hspace{0,05 cm} par pas de 0,25, les valeurs des paramètres $\varepsilon$ et $\sigma$ sont donnés dans le tableau ci-contre 


\begin{table}[h]
  \begin{flushright}
  \begin{tabular}{||l c r||}
    \hline
    \hline
    Atomes & $\varepsilon(kcla/mole)$ & $\sigma($\AA$)$ \\
    \hline
    \hline
    He & 0,12 & 2,6 \\
    Ne & 0,07 & 2,8 \\
    Ar & 0,24 & 3,4 \\ 
    Kr & 0,35 & 3,8 \\
    Xe & 0,45 & 4,0 \\
    \hline
    \hline
  \end{tabular}
  \end{flushright}
\end{table}

\noindent
Enregistre dans des fichiers (un fichier par élément gaz rare) les données dans le but de reproduirele graphe de l'énergie potentielle en foction de la position $R$. \\
Chaque fichier doit contenir deux colonnes correspondant aux valeurs succesives de $R$ et $E(R)$

\subsubsection{\underline{Exercice 2}}

\noindent
L'énergie potentielle d'une particule est donnée par la relation $ E = \frac{1}{r^3} - \frac{1}{r^2} $ .

\begin{enumerate}
  \item[(a)] Écrire un programme avec une boucle \verb|DO| permettant de calculer $E$ pour des valeurs de $r$ entières comprise entre 1 et 10. Le résultat sera enregistré dans un fichier.
  \item[(b)] Idem avec $r$ variant de 0,5 et 5 avec un pas de 0,5.
  \item[(c)] Idem avec $r$ variant de 0,7 à $X$ avec un pas de 0,3. La valeur de $X$ devrait être donnée à l'exécution.
\end{enumerate}

\subsubsection{\underline{Problème}}

\begin{enumerate}
  \item Écrire une \verb|function| externe dénomée \verb|fact(N)| permettant de calculer le factoriel d'un nombre entier positif $N$ . On rappelle que 0! = 1 .
  \item On veut calculer $e^x$ à l'aide d'un développemnt limité : 
    $$ e^x = 1 + \frac{x}{1!} + \frac{x^2}{2!} + \cdots + \frac{x^i}{i!} = a_0 + a_1 + a_2 + \cdots a_i $$ . \\
    On décide d'arrêter le développment lorsque le rapport du terme $a_i$ d'ordre $i$ sur la somme de tous les précédents est inférieur à une valeur donnée $\varepsilon $ . \\ 
    
    \vspace{0,125 cm}
    \noindent
    En se servant de la \verb|function| externe \verb|fact(N)| calculant les factoriels, écrire une \verb|function| externe \verb|exp(x)| approchant l'exponentiel d'une valeur réelle $x$ . \\
    On remarquera que chaque terme de cette somme peut être obtenu à partir du précédent par :
    $$ a_k = a_{k - 1} \times \frac{x}{k} $$ . 
  \item   
\end{enumerate}
\end{document}